\documentclass[12pt]{article}
\usepackage[margin=0.9in]{geometry}
\usepackage[utf8]{inputenc}
\usepackage[T1]{fontenc}
\usepackage{array}
\usepackage{tabularx}
\usepackage{xltabular}
\usepackage{booktabs}
\usepackage[dvipsnames]{xcolor}
\usepackage{hyperref}
\usepackage{enumitem}
\usepackage{parskip}
\usepackage{setspace}
\onehalfspacing
\definecolor{navyBlue}{HTML}{003087}
\definecolor{brickRed}{HTML}{B22222}
\definecolor{tableHeader}{HTML}{003087}
\definecolor{altRow}{HTML}{F0F4FF}
\newcolumntype{P}[1]{>{\raggedright\arraybackslash}p{#1}}
\newcolumntype{L}{>{\raggedright\arraybackslash}X}
\renewcommand{\arraystretch}{1.16}
\hypersetup{colorlinks=true,linkcolor=navyBlue,urlcolor=brickRed}
\setlist[itemize]{leftmargin=*,itemsep=2pt,topsep=3pt}

\begin{document}

\begin{center}
{\Large\bfseries\color{navyBlue} ED452B Part 2 Excellence Evidence Addendum}\par
{\large Minecraft Education Coding FUNdamentals}\par
{\normalsize Piter Garcia \textperiodcentered\ May 2, 2026}\par
\end{center}

\section*{Purpose}

This addendum strengthens the Part 2 Revised Innovative Unit Plan against the rubric's Level 4 language. The rubric asks for exceptionally illustrative examples from analyses of student data, community interactions, and/or professional collaborations to contextualize and rationalize unit design. This addendum shows how the unit design is supported by three connected evidence streams: ED452B course discussion transcripts, TeachingPlacement Minecraft classroom evidence, and the EDE498C Minecraft innovative-unit module.

\section*{Privacy and Evidence Use Note}

This public addendum summarizes evidence patterns rather than publishing raw student records, raw IEP documentation, or identifiable student work. When transcript evidence is used, it is paraphrased and de-identified. The purpose is to preserve the instructional analysis while protecting students, classmates, instructors, and placement partners.

\section{Cross-Course Evidence: EDE498C as a Parallel Minecraft Innovative Unit}

The EDE498C repository provides important professional-context evidence because it independently documents the Minecraft unit as a course-level innovative-unit module. Its public README describes the repository as a sanitized course portfolio containing assignment materials, coaching workflow resources, and a dedicated Minecraft innovative-unit module for Grade 4--5 IGNITE. It also states that the active innovative unit is anchored in a Grade 4 Minecraft Education lesson sequence for Grade 4--5 IGNITE, including basic moves and entry routines, sequencing to reach a goal, loops/debugging/revision, and efficient solutions/reflection.

This matters for ED452B because the Minecraft unit is not a one-off assignment artifact. It is a coherent cross-course instructional design problem that appears in both ED452B and EDE498C. ED452B emphasizes inclusive instruction, accommodations, assessment, target students, and unit revision. EDE498C emphasizes integrating computer science standards across the curriculum and developing professional practice. Together, the two courses support the same core claim: Minecraft Education can be a rigorous computer science environment when it is planned through inclusive routines, standards alignment, assessment validity, and classroom evidence.

The EDE498C Minecraft hub identifies a canonical innovative-unit document, a formal Warner lesson, a lesson index, a source map, and a module status file. It also names the Grade 4 formal lesson arc: Minecraft Bridge, Minecraft Tree House, Final Project Physical Build, and Project Presentation. This gives the ED452B report a stronger professional collaboration and curricular-context rationale: the ED452B inclusive unit is part of a broader curriculum-development effort, not an isolated lesson plan.

\section{Transcript Evidence Matrix}

\begin{xltabular}{\textwidth}{|P{1.45in}|P{1.9in}|P{1.85in}|L|}
\hline
\rowcolor{tableHeader}\color{white}\bfseries Evidence Pattern & \color{white}\bfseries What the Transcript Evidence Showed & \color{white}\bfseries Interpretation & \color{white}\bfseries Unit Design Response \\
\hline
Game-based access changed participation & In the March 5 ED452B check-in, the candidate described students who had been viewed as behaviorally challenging becoming engaged, seated, productive, and interested in helping peers when game-based and coding-style tasks matched their interests. & Engagement improved when instruction entered students' world and treated their digital/game knowledge as an asset. This supports Minecraft as an access point, not as a reward or gimmick. & The unit uses Minecraft Education as the core instructional environment and structures advanced students as Expert Helpers who ask questions, point to resources, and avoid taking over peers' keyboards. \\
\hline
\rowcolor{altRow} Supports preserved expectations & In the March 12 inclusive assessment discussion, classmates and instructor discussed manipulatives, teacher-guided steps, and accommodations as tools that allow students to show understanding without lowering the goal. & Support should be documented as access, not hidden as rescue. A student can use a tool and still demonstrate real understanding. & The assessment plan separates startup access, goal understanding, planning, code accuracy, debugging, and explanation so the teacher can identify what was independent and what was scaffolded. \\
\hline
Complex tasks can produce shutdown & In the March 12 meeting, the candidate described students becoming blocked when a task looks new or complex and explained that modeling a similar task in another context can help students begin without giving them the answer. & Initial refusal, random clicking, or silence may indicate cognitive load and unclear entry points rather than inability. & Lessons include teacher think-alouds, Stop--Read--Predict, Plan Before You Code, goal restatement, and one-change debugging routines. \\
\hline
\rowcolor{altRow} Assessment should be multi-source and strength-based & In the March 12 workshop, the instructor affirmed the candidate's presentation about functional assessment, especially the idea that assessment should identify strengths, supports, and conditions for success, not only struggles. & Completion alone is too narrow for inclusive computer science assessment. Student learning should be triangulated across products, process, talk, and supports. & The unit collects code artifacts, planning sheets, debugging logs, teacher conferences, oral explanations, role observations, and exit tickets. \\
\hline
Alternative products reveal hidden competence & In the March 12 workshop, a classmate described a learner whose strengths were visible through a three-dimensional model even when traditional writing was difficult. & Competence may be format-dependent. In CS, a student's strongest evidence may be oral, visual, digital, spatial, or collaborative. & Students can show learning through Agent code, screenshots, oral explanation, partner teach-back, written reflection, or teacher conference evidence. \\
\hline
\rowcolor{altRow} Relationship and safety shaped access & In the March 19 check-in, the candidate described using reassurance and relational presence when a student with significant support needs was uncomfortable. & Access is emotional and relational as well as technical. Predictable adult response supports academic participation. & The unit uses predictable agendas, calm stuck routines, quiet help signals, and non-stigmatizing support protocols. \\
\hline
Whole-class reflection avoids stigma & In the March 19 breakout and debrief, the candidate and classmates developed the idea of whole-class reflection on supports so students with formal plans are not isolated. & Whole-class reflection makes supports normal, reveals hidden needs, and helps students build metacognition. & Reflection prompts ask all students which supports helped, what was hard, and what they would try next. Support use is treated as productive learning behavior. \\
\hline
\rowcolor{altRow} Benefit versus need requires data & In the March 19 debrief, the instructor distinguished between supports that benefit many students and accommodations that a particular student needs, emphasizing evidence over automatic carryover. & IEP-related decisions should be tied to current performance, student self-report, and evidence of independence with/without supports. & Target-student sections identify evidence to collect, interpretation, feedback, future support, and possible fading of scaffolds. \\
\hline
Scaffolds should build independence & In the March 26 discussion, the class connected modifications, team decision-making, independence, and Vygotsky's Zone of Proximal Development. & Supports should keep students in productive challenge: not so hard that they shut down, not so easy that they stop growing. & The unit uses hints-before-rescue, one-change debugging, teacher checkpoints, role clarity, and extension tasks to preserve rigor while supporting access. \\
\hline
\end{xltabular}

\section{Cross-Course Design Alignment Matrix}

\begin{xltabular}{\textwidth}{|P{1.45in}|P{2.0in}|P{2.0in}|L|}
\hline
\rowcolor{tableHeader}\color{white}\bfseries Evidence Source & \color{white}\bfseries What It Adds & \color{white}\bfseries Why It Matters & \color{white}\bfseries ED452B Design Implication \\
\hline
ED452B rubric and course discussions & Inclusion, target-student analysis, IEP/accommodation logic, assessment validity, and unit revision. & Establishes the inclusive pedagogy and assessment frame required for Part 2. & The report foregrounds access barriers, assessment interpretation, target students, and future supports. \\
\hline
\rowcolor{altRow} ED452B transcripts & First-person professional reflection, peer discussion, instructor feedback, and repeated evidence about accommodations, functional assessment, student confidence, and data-informed support. & Provides exactly the kind of professional collaboration and evidence analysis needed for Level 4. & The final unit explains not only what supports are used but why those supports were chosen. \\
\hline
TeachingPlacement Minecraft evidence & Classroom grounding for routines, friction points, student engagement, Minecraft setup, NPC reading, debugging, and teacher support moves. & Prevents the unit from sounding theoretical only; ties it to actual Grade 4 IGNITE implementation conditions. & Lessons include click paths, NPC read-aloud, task-goal restatement, planning sheets, and debugging conferences. \\
\hline
\rowcolor{altRow} EDE498C Minecraft module & Parallel innovative-unit documentation, CS integration frame, formal Warner lesson arc, source map, and professional-practice framing. & Shows that the Minecraft unit is a sustained cross-course curriculum project, not a single assignment response. & The ED452B unit can claim stronger curricular coherence, cross-course professional development, and standards-aligned CS integration. \\
\hline
\end{xltabular}

\section{Revised Excellence-Level Rationale Paragraph}

The following paragraph may be inserted into the main Part 2 report's Unit Context and Rationale section:

\begin{quote}
The rationale for this unit is substantiated by more than one evidence stream. ED452B transcripts show repeated professional discussion around inclusive assessment, accommodations, student independence, functional assessment, and the need to distinguish support from lowered expectations. TeachingPlacement evidence shows that Minecraft tasks create both high engagement and predictable access barriers, including world entry, reading NPC directions, interpreting multi-step goals, keyboard/mouse control, and debugging. EDE498C provides a parallel cross-course Minecraft innovative-unit module, confirming that this lesson sequence is part of a broader computer science integration project rather than a stand-alone activity. Together, these sources justify the unit's central design decision: assess computational thinking through multiple evidence types while explicitly teaching the access routines that make that thinking visible.
\end{quote}

\section{Revised Excellence-Level Systematic Analysis Paragraph}

The following paragraph may be inserted into the main Part 2 report's Systematic Analysis of Student Learning section:

\begin{quote}
The student-learning analysis is organized around a pattern that emerged across classroom evidence and ED452B transcript analysis: students' visible performance depends on the match between the task format, support structure, and assessment mode. In Minecraft, a student may fail to complete a challenge because of a coding misconception, but the same surface failure can also come from login friction, NPC reading load, unclear success criteria, keyboard/mouse demands, partner-role confusion, or frustration after the Agent behaves unexpectedly. For that reason, the unit does not treat completion as the only evidence of learning. It triangulates code artifacts, planning sheets, debugging conversations, oral explanations, exit tickets, and support-use observations to identify what the student understands, what support revealed that understanding, and what future instruction should fade or maintain.
\end{quote}

\section{Why This Moves the Report Toward Level 4}

This evidence strengthens the report in four ways:

\begin{itemize}
\item It gives exceptionally illustrative examples from professional collaboration by using ED452B transcript discussions rather than generic claims.
\item It uses student-data logic by analyzing access, planning, debugging, explanation, and support-use patterns separately.
\item It contextualizes the unit through community and placement evidence from Pine Brook IGNITE and TeachingPlacement Minecraft lessons.
\item It substantiates the rationale with theory: UDL, MTSS, Karten's inclusion framework, functional assessment, and ZPD-informed scaffolding.
\end{itemize}

\section{Submission Use}

This addendum can be submitted as an appendix item, or the two revised paragraphs above can be copied directly into the main Part 2 report. The strongest submission would do both: keep this addendum as appendix evidence and integrate the shorter paragraphs into the main report so the Level 4 rationale is visible before the appendix.

\end{document}
